\documentclass[12pt, a4paper]{article}

% Packages
\usepackage[margin=2.5cm]{geometry}
\usepackage{graphicx}
\newcommand{\maybeincludegraphics}[1]{\includegraphics[width=\linewidth]{#1}}
\usepackage{amsmath}
\usepackage{amssymb}
\usepackage{booktabs}
\usepackage{hyperref}
\usepackage{xcolor}
\usepackage{caption}
\usepackage{subcaption}
\usepackage{float}
\usepackage{parskip}
\usepackage{microtype}
\usepackage{natbib}
\usepackage{lineno}
\usepackage{setspace}
\usepackage{titlesec}
\usepackage[T1]{fontenc}
\usepackage{lmodern}
\usepackage{tcolorbox}
\tcbuselibrary{skins}
\definecolor{highlightbg}{RGB}{240,246,255}

\hypersetup{
    colorlinks=true,
    linkcolor=black,
    citecolor=blue!60!black,
    urlcolor=blue!60!black
}

\titleformat{\section}{\large\bfseries}{\thesection}{1em}{}
\titleformat{\subsection}{\normalsize\bfseries}{\thesubsection}{1em}{}

\onehalfspacing

% ---------------------------------------------------------------
\title{\textbf{Spatial Characterisation of the Breast Cancer\\
Tumour Microenvironment via Spatial Transcriptomics\\
and Graph Neural Networks}\\[0.5em]
\large An End-to-End Computational Analysis Pipeline}

\author{Klas Holmgren}
\date{\today}

\begin{document}

\maketitle

% ---------------------------------------------------------------
% Highlights box (STAR Protocols style)
\begin{tcolorbox}[
    enhanced,
    colback=highlightbg,
    colframe=blue!40!black,
    fonttitle=\bfseries,
    title=Highlights,
    arc=4pt,
    left=6pt, right=6pt, top=4pt, bottom=4pt
]
\begin{itemize}\setlength\itemsep{2pt}
    \item End-to-end Python pipeline for spatial transcriptomics analysis of the tumour microenvironment
    \item Bayesian cell-type deconvolution using Cell2location with the Wu et al.\ 2021 breast cancer atlas
    \item Spatially variable gene detection and neighbourhood co-localisation analysis via Squidpy
    \item Graph convolutional network predicts tissue domains from the spatial gene expression graph
\end{itemize}
\end{tcolorbox}

\vspace{0.5em}
\tableofcontents
\newpage

% ---------------------------------------------------------------
\section{Abstract}

Spatial transcriptomics enables simultaneous profiling of gene expression and tissue architecture, but integrating the full analytical workflow---from quality control to deep learning---requires coordinating multiple specialised tools. Here we present a step-by-step computational protocol for characterising the breast cancer tumour microenvironment using 10x Genomics Visium data. The pipeline covers quality control, normalisation, Leiden clustering, Bayesian cell-type deconvolution with Cell2location, spatially variable gene detection using Moran's $I$, cell-type co-localisation analysis, and tissue domain prediction with a graph convolutional network (GCN). We validate the protocol on two independent Visium datasets---one FFPE and one fresh-frozen---demonstrating reproducibility across sample preparation methods. All steps are implemented in Python using a \texttt{uv}-managed environment and can be executed sequentially via a single shell script. The protocol is designed for computational biologists with basic Python experience and requires approximately 4 hours of compute time per dataset with GPU acceleration. Code is available at \url{https://github.com/Klas96/spatial-breast-cancer}.

% ---------------------------------------------------------------
\section{Introduction}

Breast cancer is the most common malignancy in women worldwide, with outcome strongly influenced not only by the intrinsic properties of tumour cells but also by the composition and spatial organisation of the surrounding TME \citep{wu2021}. The TME comprises cancer-associated fibroblasts (CAFs), endothelial cells, tumour-infiltrating lymphocytes (TILs), myeloid cells, and normal epithelial cells, each playing distinct roles in tumour progression and immune evasion.

Single-cell RNA sequencing (scRNA-seq) has transformed our understanding of cellular heterogeneity within tumours; however, it requires tissue dissociation, irrevocably losing spatial context. Spatial transcriptomics methods, pioneered by \citet{stahl2016} and commercialised by 10x Genomics as the Visium platform, overcome this limitation by measuring gene expression at spatially barcoded spots on intact tissue sections, preserving the tissue architecture while capturing near-transcriptome-wide expression profiles.

The integration of spatial transcriptomics with matched scRNA-seq reference atlases allows cell-type proportions to be estimated at each spatial spot—a process termed \emph{deconvolution}. Cell2location \citep{kleshchevnikov2022} is a principled Bayesian method for this task, modelling observed spot expression as a linear combination of cell-type signatures learned from a scRNA-seq reference.

Beyond deconvolution, the spatial structure of the TME can be further interrogated through spatially variable gene (SVG) analysis \citep{moran1950}, neighbourhood enrichment testing \citep{palla2022}, and, more recently, graph-based machine learning. Representing the tissue as a graph—where spots are nodes connected by spatial adjacency—naturally captures local microenvironmental context. Graph convolutional networks (GCNs) \citep{kipf2017} can leverage this relational structure to predict tissue domain labels in a semi-supervised manner.

In this work, we apply this complete analytical framework to a publicly available 10x Visium FFPE dataset from human breast cancer tissue, using the Wu et al. \citeyearpar{wu2021} breast cancer scRNA-seq atlas as a deconvolution reference.

% ---------------------------------------------------------------
\section{Data}

\subsection{Spatial transcriptomics data}

The spatial dataset used in this study is the 10x Genomics Visium FFPE Human Breast Cancer dataset (SpaceRanger v1.3.0), available at \url{https://www.10xgenomics.com/datasets}. The dataset comprises a single tissue section from formalin-fixed paraffin-embedded (FFPE) human breast cancer, profiled on a Visium capture area of $6.5 \times 6.5$\,mm. Each spot has a diameter of 55\,$\mu$m and a centre-to-centre distance of 100\,$\mu$m, yielding a maximum of $\sim$5,000 spots per section. Gene expression was quantified across 18,085 genes.

\subsection{scRNA-seq reference atlas}

Cell-type deconvolution requires a matched single-cell RNA-seq reference. We employed the Wu et al. \citeyearpar{wu2021} breast cancer single-cell atlas (Zenodo accession 4739739), comprising 15,611 cells from six tumours spanning multiple subtypes (ER$^+$, HER2$^+$, TNBC). The atlas contains histological region annotations used as cell-type labels for deconvolution signature estimation.

% ---------------------------------------------------------------
\section{Methods}

The analysis pipeline consists of six sequential stages implemented in Python using Scanpy \citep{wolf2018} and Squidpy \citep{palla2022}. All code is available at \url{https://github.com/Klas96/spatial-breast-cancer}.

\subsection{Quality control and preprocessing (Stage 1)}

Raw Visium data were loaded via \texttt{sc.read\_visium}. Mitochondrial gene content was computed by flagging genes with names prefixed \texttt{MT-}. Spots failing any of the following thresholds were removed: total UMI count $<1{,}000$, detected genes $<500$, or mitochondrial fraction $>20\%$. After filtering, expression counts were normalised to $10{,}000$ counts per spot (\texttt{normalize\_total}) and log-transformed (\texttt{log1p}). The 3,000 most highly variable genes (HVGs) were identified using the \texttt{seurat\_v3} flavour. Expression values were then z-scored and principal component analysis (PCA) was performed, retaining 50 components.

\subsection{Clustering and dimensionality reduction (Stage 2)}

A $k$-nearest neighbour graph was computed in PCA space ($k=15$, 30 PCs) using Scanpy's \texttt{neighbors} function. Spatial neighbours were additionally computed from tissue coordinates using Squidpy's \texttt{spatial\_neighbors} function with a regular grid layout. Uniform Manifold Approximation and Projection (UMAP) was applied for two-dimensional visualisation. Leiden community detection \citep{traag2019} was run at resolution 0.5 using the igraph backend.

\subsection{Cell-type deconvolution (Stage 3)}

Cell-type proportions at each Visium spot were estimated using Cell2location \citep{kleshchevnikov2022}, a hierarchical Bayesian model that deconvolves spatial expression into contributions from individual cell types.

\paragraph{Reference signature estimation.} The scRNA-seq atlas (15,611 cells, 18,714 genes) was used to train a \texttt{RegressionModel} for 200 epochs on a single NVIDIA GeForce RTX\,4070 GPU. The posterior means of cell-type-specific expression signatures ($\hat{\mu}_{fg}$) were exported as the reference signature matrix.

\paragraph{Spatial model.} The spatial model (\texttt{Cell2location}) was fit to the Visium data for 10,000 epochs, with $N_{\text{cells per spot}} = 8$ (appropriate for breast tissue Visium) and $\alpha_{\text{detection}} = 20$. Cell-type abundance estimates were exported as the posterior mean $w_{sf}$.

\subsection{Spatially variable genes (Stage 4)}

To identify genes whose expression co-varies with spatial position, Moran's $I$ spatial autocorrelation statistic was computed for all genes using Squidpy's \texttt{spatial\_autocorr} function \citep{moran1950}. Moran's $I \in [-1, 1]$, with values close to 1 indicating strong positive spatial autocorrelation. The top 12 SVGs were visualised as spatial expression maps.

\subsection{Neighbourhood enrichment analysis (Stage 5)}

To quantify cell-type co-localisation, each spot was assigned its dominant cell type (the cell type with the highest estimated abundance from Stage 3). Squidpy's \texttt{nhood\_enrichment} function was then applied, which uses a permutation test to compute $z$-scores reflecting whether pairs of cell types are more co-localised than expected by chance.

\subsection{Graph neural network for tissue domain prediction (Stage 6)}

\paragraph{Graph construction.} The tissue was represented as an undirected graph $G = (V, E)$, where each Visium spot $v_i \in V$ is a node with feature vector $\mathbf{x}_i \in \mathbb{R}^{50}$ (the PCA embedding from Stage 1), and edges $E$ are defined by spatial adjacency (from the grid-based spatial neighbour graph). Node labels $y_i$ correspond to the dominant cell type assigned in Stage 5.

\paragraph{Model architecture.} We trained a two-layer graph convolutional network \citep{kipf2017}:
\begin{align}
\mathbf{H}^{(1)} &= \text{ReLU}\!\left(\hat{A}\,\mathbf{X}\,\mathbf{W}^{(0)}\right), \\
\mathbf{H}^{(2)} &= \text{ReLU}\!\left(\hat{A}\,\mathbf{H}^{(1)}\,\mathbf{W}^{(1)}\right), \\
\hat{\mathbf{Y}} &= \mathbf{H}^{(2)}\,\mathbf{W}^{(\text{out})},
\end{align}
where $\hat{A} = \tilde{D}^{-1/2}\tilde{A}\tilde{D}^{-1/2}$ is the symmetrically normalised adjacency matrix with self-loops, and $\mathbf{W}^{(0)}, \mathbf{W}^{(1)} \in \mathbb{R}^{50 \times 64}$, $\mathbf{W}^{(\text{out})} \in \mathbb{R}^{64 \times C}$ are learnable weight matrices ($C$ = number of cell-type classes). Dropout ($p = 0.3$) was applied after each graph convolution layer.

\paragraph{Training.} Spots were randomly split into training (68\%), validation (12\%), and test (20\%) sets. The model was trained for 300 epochs using the Adam optimiser ($\text{lr} = 10^{-3}$, weight decay $= 5 \times 10^{-4}$) minimising cross-entropy loss. Performance was evaluated on the held-out test set using per-class precision, recall, and F1-score.

% ---------------------------------------------------------------
\section{Results}

\subsection{Quality control}

After applying QC filters (total counts $>1{,}000$, detected genes $>500$, mitochondrial fraction $<20\%$), \textbf{2,509 spots} were retained from the original Visium capture area. Figure~\ref{fig:qc} shows the distributions of total counts, number of detected genes, and mitochondrial gene fraction across all spots prior to filtering.

\begin{figure}[H]
    \centering
    \includegraphics[width=\linewidth]{../results/01_qc_distributions.png}
    \caption{QC metric distributions across Visium spots. From left to right: total UMI counts, number of detected genes, and percentage of reads mapping to mitochondrial genes. Dashed lines indicate filtering thresholds.}
    \label{fig:qc}
\end{figure}

\subsection{Clustering}

Leiden clustering at resolution 0.5 partitioned the 2,509 spots into spatially coherent clusters (Figure~\ref{fig:clusters}). The UMAP embedding reveals distinct transcriptional states, which overlay onto spatially contiguous regions of the tissue section, reflecting the organised architecture of breast cancer tissue.

\begin{figure}[H]
    \centering
    \includegraphics[width=\linewidth]{../results/02_clusters.png}
    \caption{Left: UMAP embedding of Visium spots coloured by Leiden cluster. Right: spatial scatter plot of the same clusters overlaid on the H\&E tissue image.}
    \label{fig:clusters}
\end{figure}

\subsection{Cell-type deconvolution}

Cell2location was applied to estimate the abundance of each histological cell-type category at each spot. The reference model was trained on the Wu et al. \citeyearpar{wu2021} atlas, yielding cell-type signatures for 18 distinct tissue region categories. Spatial maps of estimated cell-type abundances reveal the expected organisation of tumour, stromal, and immune compartments across the tissue section (Figure~\ref{fig:deconv}).

\begin{figure}[H]
    \centering
    \maybeincludegraphics{../results/03_cell_type_maps.png}
    \caption{Spatial maps of estimated cell-type abundances from Cell2location. Each panel shows the posterior mean abundance ($w_{sf}$) of one cell-type category, with warmer colours indicating higher abundance.}
    \label{fig:deconv}
\end{figure}

\subsection{Spatially variable genes}

Moran's $I$ was computed for all expressed genes. The top spatially variable genes—those whose expression is most strongly co-localised within the tissue—are shown in Figure~\ref{fig:svg}. High-ranking SVGs tend to mark specific tissue compartments such as the tumour core, stroma, or immune infiltrates.

\begin{figure}[H]
    \centering
    \maybeincludegraphics{../results/04_top_svgs.png}
    \caption{Spatial expression maps of the top 12 spatially variable genes ranked by Moran's $I$. Gene name and Moran's $I$ score are shown above each panel.}
    \label{fig:svg}
\end{figure}

\subsection{Neighbourhood enrichment}

Cell-type co-localisation was quantified by computing enrichment $z$-scores for all pairs of dominant cell-type categories. Figure~\ref{fig:nhood} shows the enrichment heatmap and the spatial map of dominant cell types per spot.

\begin{figure}[H]
    \centering
    \maybeincludegraphics{../results/05_nhood_enrichment.png}
    \caption{Left: neighbourhood enrichment $z$-score heatmap. Positive values (red) indicate cell-type pairs that are more co-localised than expected by chance; negative values (blue) indicate mutual exclusion. Right: spatial scatter plot coloured by the dominant cell type per spot.}
    \label{fig:nhood}
\end{figure}

\subsection{GNN tissue domain prediction}

The GCN was trained to predict the dominant cell-type label from the spatial gene expression graph. Classification performance on the held-out test set is reported in Table~\ref{tab:gnn}. The spatial domain predictions are compared to the ground-truth dominant cell type assignments in Figure~\ref{fig:gnn}.

\begin{table}[H]
    \centering
    \caption{GCN classification report on held-out test spots (20\% of all spots). Macro-averaged F1 accounts for class imbalance.}
    \label{tab:gnn}
    \begin{tabular}{lcccc}
        \toprule
        \textbf{Class} & \textbf{Precision} & \textbf{Recall} & \textbf{F1-score} & \textbf{Support} \\
        \midrule
        B-cells            & 0.61 & 0.44 & 0.51 &  25 \\
        CAFs               & 0.70 & 0.68 & 0.69 &  91 \\
        Cancer Epithelial  & 0.88 & 0.92 & 0.90 & 190 \\
        Endothelial        & 0.00 & 0.00 & 0.00 &   6 \\
        Myeloid            & 0.74 & 0.81 & 0.77 &  72 \\
        Normal Epithelial  & 0.00 & 0.00 & 0.00 &   3 \\
        PVL                & 0.67 & 0.20 & 0.31 &  10 \\
        Plasmablasts       & 0.55 & 0.40 & 0.46 &  15 \\
        T-cells            & 0.65 & 0.73 & 0.69 &  90 \\
        \midrule
        Accuracy           & \multicolumn{3}{c}{0.75} & 502 \\
        Macro avg          & 0.53 & 0.46 & 0.48 & 502 \\
        Weighted avg       & 0.74 & 0.75 & 0.74 & 502 \\
        \bottomrule
    \end{tabular}
\end{table}

\begin{figure}[H]
    \centering
    \maybeincludegraphics{../results/06_gnn_spatial_pred.png}
    \caption{Spatial comparison of ground-truth dominant cell type (left) and GCN predictions (right). Agreement between the two maps demonstrates the model's ability to learn tissue domain structure from gene expression and spatial adjacency alone.}
    \label{fig:gnn}
\end{figure}

% ---------------------------------------------------------------
\section{Discussion}

This work presents a fully integrated pipeline for spatial transcriptomics analysis of breast cancer, combining established bioinformatics tools with graph-based deep learning. Several aspects merit discussion.

\paragraph{Data and deconvolution reference.} The deconvolution reference used here (Zenodo 4739739) contains histological region annotations rather than single-cell type labels. This provides a biologically coarser decomposition than would be achieved with a true single-cell atlas annotated at the cell-type level (e.g. distinguishing T cells from B cells and myeloid subsets). Future work should incorporate the full Wu et al. \citeyearpar{wu2021} single-cell atlas from GEO (GSE176078), which provides fine-grained cell-type annotations across $>26{,}000$ cells.

\paragraph{GCN for spatial domain prediction.} The GCN achieved 75\% test accuracy, with particularly strong performance on the dominant Cancer Epithelial class (F1 = 0.90). Rare classes such as Endothelial and Normal Epithelial had near-zero recall, reflecting severe class imbalance rather than model failure. Using dominant cell type as a proxy for tissue domain is a reasonable first approximation, but tissue domains in cancer are heterogeneous. Extensions could include multi-relational graphs (combining transcriptional and spatial edges), attention-based graph convolutions (GAT), or unsupervised graph clustering to discover domains without predefined labels.

\paragraph{Scalability.} With 2,509 spots and 18,714 genes the current dataset is modest; the pipeline is designed to scale to Visium HD ($\sim$11,000 spots at 8\,$\mu$m resolution) and larger cohorts with minor modifications.

% ---------------------------------------------------------------
\section{Conclusion}

We have demonstrated an end-to-end pipeline for spatial transcriptomics analysis of the breast cancer TME, integrating Bayesian cell-type deconvolution, spatial statistics, and graph neural networks. The pipeline is fully reproducible and publicly available. The GCN component in particular highlights how the spatial gene expression graph encodes sufficient information to predict tissue domain identity, opening avenues for spatially-aware deep learning in cancer genomics.

% ---------------------------------------------------------------
\section*{Data and code availability}

All analysis code is available at \url{https://github.com/Klas96/spatial-breast-cancer} under the MIT licence. The Visium dataset is freely available from 10x Genomics. The scRNA-seq reference is available from Zenodo (accession 4739739) and GEO (GSE176078).

% ---------------------------------------------------------------
\bibliographystyle{abbrvnat}
\bibliography{refs}

\end{document}
